\DocumentMetadata{
  pdfversion=2.0,
  pdfstandard=ua-2,
  lang=en-US,
  tagging=on,
  tagging-setup={math/setup=mathml-SE,tabsorder=structure}
}
\documentclass[11pt,letterpaper]{article}
\usepackage[margin=0.68in,includefoot]{geometry}
\usepackage{newcomputermodern}
\usepackage{amsmath,amscd,mathtools}
\usepackage{tikz,adjustbox}
\providecommand{\square}{\mdlgwhtsquare}
\usepackage[hidelinks]{hyperref}
\hypersetup{
  pdftitle={Numerical Analysis PhD Exam, January 2022},
  pdfauthor={University of Florida Department of Mathematics},
  pdfsubject={Graduate mathematics examination},
  pdfdisplaydoctitle=true,
  bookmarksopen=true
}
\setcounter{secnumdepth}{0}
\setlength{\parindent}{0pt}
\setlength{\parskip}{0.28em}
\setlength{\emergencystretch}{3em}
\setlength{\leftmargini}{2.15em}
\setlength{\leftmarginii}{2.65em}
\setlength{\leftmarginiii}{3em}
\pagestyle{empty}
\raggedbottom
\allowdisplaybreaks
\NewTaggingSocketPlug{block/list/label}{examproblem}{%
  \tagstructbegin{tag=Lbl}\tagstructend%
  \tagstructbegin{tag=\LBody}%
  \tagstructbegin{tag=\ProblemHeadingTag,title-o={\ProblemHeadingText}}%
  \tagmcbegin{tag=\ProblemHeadingTag,actualtext={\ProblemHeadingText}}%
  #2%
  \tagmcend%
  \tagstructend%
}
\newenvironment{examproblems}[2]{%
  \def\ProblemHeadingTag{#2}%
  \AssignTaggingSocketPlug{block/list/label}{examproblem}%
  \begin{enumerate}%
  \setcounter{enumi}{#1}%
  \renewcommand{\labelenumi}{\textbf{\theenumi.}}%
  \setlength{\topsep}{0.35em}%
  \setlength{\partopsep}{0pt}%
  \setlength{\itemsep}{0.48em}%
  \setlength{\parsep}{0.18em}%
}{\end{enumerate}}
\newenvironment{examsubparts}{%
  \AssignTaggingSocketPlug{block/list/label}{default}%
  \begin{enumerate}%
  \setlength{\topsep}{0.18em}%
  \setlength{\partopsep}{0pt}%
  \setlength{\itemsep}{0.16em}%
  \setlength{\parsep}{0.1em}%
}{\end{enumerate}}
\newenvironment{examinstructions}{%
  \par\begingroup\itshape
}{%
  \par\endgroup\vspace{0.18em}
}
\makeatletter
\renewcommand\subsection{\@startsection{subsection}{2}{0pt}%
  {-1.25ex plus -.25ex minus -.1ex}%
  {0.55ex plus .1ex}%
  {\normalfont\large\bfseries}}
\renewcommand{\@maketitle}{%
  \newpage\null\vskip -1em%
  {\footnotesize\bfseries
    \href{https://www.ufl.edu/}{University of Florida}, \href{https://math.ufl.edu/}{Department of Mathematics}\par}%
  \vskip -0.5em%
  \tagstructbegin{tag=H1,title={Numerical Analysis PhD Exam, January 2022}}%
  \tagpdfparaOff%
  \tagmcbegin{tag=H1}%
    {\LARGE\bfseries\href{https://gma.math.ufl.edu/exams/numerical-analysis-phd/}{Numerical Analysis PhD Exam}, January 2022\par}%
  \tagmcend%
  \tagpdfparaOn%
  \tagstructend%
  \vskip 0.25em%
  \tagmcbegin{artifact}%
    \hrule height 1pt%
  \tagmcend%
  \vskip 1em%
}
\makeatother
\title{Numerical Analysis PhD Exam, January 2022}
\author{}
\date{}
\begin{document}
\maketitle
\thispagestyle{empty}
\begin{examinstructions}
Do 4 (four) of the first 5 (1–5) and 4 (four) of the last 5 problems (6–10).
\end{examinstructions}
\begin{examproblems}{0}{H2}
\def\ProblemHeadingText{Problem 1.}
\item
Let \(A \in \mathbb{C}^{m\times n}\) and \(B \in \mathbb{C}^{n\times p}\). Prove or provide a counterexample to the following statements. Justify each nontrivial step.
\begin{examsubparts}
\item[\textnormal{(a)}]
The Frobenius norm satisfies \(\lVert AB\rVert_F \leq \lVert A\rVert_F\lVert B\rVert_F\).
\item[\textnormal{(b)}]
\(\lVert A\rVert_2 \leq \lVert A\rVert_F\), where \(\lVert A\rVert_F\) is the Frobenius norm of~\(A\).
\end{examsubparts}
\def\ProblemHeadingText{Problem 2.}
\item
Define the matrices~\(A\) and~\(B\) by
\[
A=\begin{pmatrix}1&2&3\\2&4&6\\1&2&3\\2&4&6\end{pmatrix},\qquad B=\begin{pmatrix}1&1&0\\0&1&3\\0&1&3\\2&0&0\end{pmatrix}.
\]
\begin{examsubparts}
\item[\textnormal{(a)}]
Find an economy singular value decomposition of~\(A\).
\item[\textnormal{(b)}]
Find an economy QR decomposition of~\(B\).
\end{examsubparts}
\def\ProblemHeadingText{Problem 3.}
\item
Let \(\lVert\cdot\rVert\) be a subordinate (induced) matrix norm.
\begin{examsubparts}
\item[\textnormal{(a)}]
If~\(E\) is \(n\times n\) with \(\lVert E\rVert<1\), then show \(I+E\) is nonsingular and
\[
\lVert(I+E)^{-1}\rVert\leq \frac{1}{1-\lVert E\rVert}.
\]
\item[\textnormal{(b)}]
If~\(A\) is \(n\times n\) invertible and~\(E\) is \(n\times n\) with \(\lVert A^{-1}\rVert\lVert E\rVert<1\), then show \(A+E\) is nonsingular and
\[
\lVert(A+E)^{-1}\rVert\leq\frac{\lVert A^{-1}\rVert}{1-\lVert A^{-1}\rVert\lVert E\rVert}.
\]
\end{examsubparts}
\def\ProblemHeadingText{Problem 4.}
\item
Let \(P\in\mathbb{C}^{m\times m}\) be a projector. Show that \(\lVert P\rVert_2=1\) if and only if~\(P\) is an orthogonal projector.
\def\ProblemHeadingText{Problem 5.}
\item
Suppose~\(A\) is Hermitian positive definite.
\begin{examsubparts}
\item[\textnormal{(a)}]
Prove that each principal submatrix of~\(A\) is Hermitian positive definite.
\item[\textnormal{(b)}]
Prove that an element of~\(A\) with largest magnitude lies on the diagonal.
\item[\textnormal{(c)}]
Prove that~\(A\) has a Cholesky decomposition.
\end{examsubparts}
\def\ProblemHeadingText{Problem 6.}
\item
Consider \(I=\int_{-1}^{1}|x|\,dx=1\). Let \(T(h)\) be the composite trapezoidal rule approximation to~\(I\), for \(h=2/n\), \(n\in\mathbb{N}\).
\begin{examsubparts}
\item[\textnormal{(a)}]
Show \(T(2/n)=1\) whenever~\(n\) is even.
\item[\textnormal{(b)}]
Find \(T(2/n)\) when~\(n\) is odd, and determine the order of convergence in terms of~\(h\).
\end{examsubparts}
\def\ProblemHeadingText{Problem 7.}
\item
This problem has two parts.
\begin{examsubparts}
\item[\textnormal{(a)}]
Let \(x_1=x_0+h\) and \(x_2=x_0+2h\). Determine (explicitly find) an \(O(h^2)\) difference approximation to \(f'(x_0)\) based on \(f(x_0)\), \(f(x_1)\) and \(f(x_2)\).
\item[\textnormal{(b)}]
Prove your solution to (a) is \(O(h^2)\)
\end{examsubparts}
\def\ProblemHeadingText{Problem 8.}
\item
This problem has two parts.
\begin{examsubparts}
\item[\textnormal{(a)}]
Consider the inner product on \(C(0,2)\) given by \((f,g)=\int_0^2 f(t)g(t)\,dt\). Find three orthonormal polynomials \(\phi_0,\phi_1,\phi_2\) on \((0,2)\) with respect to the given inner product such that the degree of \(\phi_n\) is equal to~\(n\), \(n=0,1,2\).
\item[\textnormal{(b)}]
Find the nodes \(t_1\) and \(t_2\) and weights \(w_1\) and \(w_2\) which yield the weighted Gaussian Quadrature formula
\[
\int_0^2 f(t)\,dt\approx w_1f(t_1)+w_2f(t_2)
\]
 with degree of exactness \(m=3\). You should find the nodes exactly, and may leave the weights \(w_1,w_2\) in integral form.
\end{examsubparts}
\def\ProblemHeadingText{Problem 9.}
\item
Let \(x_0,x_1,\ldots,x_n\) be \(n+1\) distinct numbers. Let \(l_j(x)\) be the associated Lagrange basis polynomials, \(j=0,\ldots,n\).
\begin{examsubparts}
\item[\textnormal{(a)}]
State the definition of \(l_j(x)\) and show that \(\{l_j(x)\}_{j=0}^n\) form a basis for \(\mathcal{P}_n\), the space of polynomials of degree at most~\(n\).
\item[\textnormal{(b)}]
Show that
\[
\sum_{j=0}^n (x-x_j)^k l_j(x)=0,\qquad \text{for all }k=1,\ldots,n.
\]
\end{examsubparts}
\def\ProblemHeadingText{Problem 10.}
\item
The Littlewood–Salem–Izumi constant \(\alpha_0\) is the unique solution \(\alpha_0\in(0,1)\) of
\[
\int_0^{3\pi/2}\frac{\cos(t)}{t^\alpha}\,dt=0.
\]
 Describe how you could use Newton’s method together with a composite Clenshaw–Curtis quadrature rule based on the \(\mathcal{P}_1\) interpolant over~\(n\) subintervals to approximate \(\alpha_0\). Be specific about how each function used in the Newton method is defined, and be specific about how the nodes and weights for the quadrature rule are defined.

\par
For reference, the first few Chebyshev polynomials are \(T_0(x)=1\), \(T_1(x)=x\), \(T_2(x)=2x^2-1\), \(T_3(x)=4x^3-3x\).
\end{examproblems}
\end{document}
